% Part: normal-modal-logic
% Chapter: axioms-systems
% Section: derived-rules

\documentclass[../../../include/open-logic-section]{subfiles}

\begin{document}

\olfileid[pt]{nml}{prf}{der}

\olsection{Regras derivadas}

Encontrar e escrever !!{derivation}s é, evidentemente, difícil,
trabalhoso e repetitivo. Por exemplo, com muita frequência queremos
passar de $!A \lif !B$ para $\Box !A \lif \Box !B$, O.s.\@, aplicar a
regra~\RK. Isso exige uma aplicação de \Nec, o registro da instância
apropriada de~\Ax{K} e, em seguida, uma aplicação de~\MP. Passar de
$!A \lif !B$ e $!B \lif !C$ para $!A \lif !C$ exige registrar a
(longa) instância tautológica
\[
(!A \lif !B) \lif ((!B \lif !C) \lif (!A \lif !C))
\]
e aplicar \MP{} duas vezes. Muitas vezes queremos substituir uma
sub-!!{formula} por uma fórmula que sabemos ser equivalente, p.ex., \iftag{prvDiamond}{$\Diamond !A$ por $\lnot\Box\lnot
  !A$}{$\Box !A$ por $\lnot\Diamond\lnot !A$}, ou $\lnot\lnot !A$ por
$!A$. Assim, em vez de escrever a !!{derivation} propriamente dita, é
mais conveniente apenas registrar por que as etapas intermediárias
são !!{derivable}. Para esse fim, reunamos alguns fatos sobre
!!{derivability}.

\begin{prop}
  Se $\Log{K} \Proves !A_1$, \dots, $\Log{K} \Proves !A_n$, e $!B$
  decorre de $!A_1$, \dots, $!A_n$ pela lógica proposicional, então
  $\Log{K} \Proves !B$.
\end{prop}

\begin{proof}
  Se $!B$ decorre de $!A_1$, \dots,~$!A_n$ pela lógica proposicional, então
  \[
  !A_1 \lif (!A_2 \lif \cdots (!A_n \lif !B)\dots)
  \]
  é uma instância tautológica. Aplicar \MP{} $n$ vezes produz uma
  !!{derivation} de~$!B$.
\end{proof}

Indicaremos o uso desta proposição por~\PL.

\begin{prop}
  Se $\Log{K} \Proves !A_1 \lif (!A_2 \lif \cdots (!A_{n-1} \lif
  !A_n)\dots)$, então $\Log{K} \Proves \Box !A_1 \lif (\Box !A_2 \lif
  \cdots (\Box !A_{n-1} \lif \Box !A_n)\dots)$.
\end{prop}

\begin{proof}
  Por indução em $n$, exatamente como na demonstração de
  \olref[nor]{prop:rk}.
\end{proof}

Indicaremos o uso desta proposição por~\RK. Ilustremos como esses
resultados ajudam a estabelecer resultados de !!{derivability} com
maior facilidade.

\begin{prop}
  $\Log{K} \Proves (\Box!A \land \Box!B) \lif \Box (!A \land !B)$
\end{prop}

\begin{proof}
  \begin{derivation}
    1. & $\Log{K} \Proves !A \lif (!B \lif (!A \land !B))$ & \Taut \\
    2. & $\Log{K} \Proves \Box!A \lif (\Box !B \lif \Box(!A \land !B)))$ & \RK, 1\\
    3. & $\Log{K} \Proves (\Box!A \land \Box!B) \lif \Box(!A \land !B)$ & \PL, 2
  \end{derivation}
\end{proof}

\begin{prop}\ollabel{prop:rewriting}
  Se $\Log{K} \Proves !A \liff !B$ e $\Log{K} \Proves
  \Subst{!C}{!A}{q}$, então $\Log{K} \Proves \Subst{!C}{B}{q}$
\end{prop}

\begin{proof}
  Exercício.
\end{proof}

\begin{prob}
  Demonstre \olref[nml][prf][der]{prop:rewriting} provando, por indução
  na complexidade de~$!C$, que, se $\Log{K} \Proves !A \liff !B$,
  então $\Log{K} \Proves \Subst{!C}{!A}{q} \liff \Subst{!C}{!B}{q}$.
\end{prob}

Esta proposição é especialmente útil quando queremos converter
$\Diamond$ em $\Box$ (ou vice-versa), ou eliminar negações duplas no
interior de !!a{formula}. Daqui em diante, marcaremos as aplicações de
\olref{prop:rewriting} por ``$!A$ por $!B$'' sempre que reescrevermos
!!a{formula}~$!C(!B)$ no lugar de~$!C(!A)$. Em outras palavras,
``$!A$ por $!B$'' abrevia:
  \begin{derivation}
    & $\Proves !C(!A)$ \\
    & $\Proves !A \liff !B$\\
    & $\Proves !C(!B)$ & por \olref{prop:rewriting}
  \end{derivation}
Por exemplo:

\begin{prop}
  $\Log{K} \Proves \lnot\Box p \lif \Diamond \lnot p$
\end{prop}

\begin{proof}
  \iftag{prvDiamond}{% Diamond é primitivo
  \begin{derivation}
    1. & $\Log{K} \Proves \Diamond \lnot p \liff \lnot\Box\lnot\lnot p$ &
    \Dual\\
    2. & $\Log{K} \Proves \lnot \Box \lnot\lnot p \lif \Diamond \lnot p$ & \PL, 1\\
    3. & $\Log{K} \Proves \lnot\Box p \lif \Diamond \lnot p$ & $p$ por $\lnot\lnot p$\\
  \end{derivation}
  }{% Diamond é definido
  \begin{derivation}
    1. & $\Log{K} \Proves \lnot \Box \lnot\lnot p \lif \Diamond \lnot p$ & \Taut\\
    2. & $\Log{K} \Proves \lnot\Box p \lif \Diamond \lnot p$ & $p$ por $\lnot\lnot p$\\
  \end{derivation}
  A fórmula na linha~$1$ é uma instância de $\lnot q \lif \lnot q$,
  na qual $\Box\lnot\lnot p$ substitui~$q$. $\Diamond\lnot p$ e
  $\lnot\Box\lnot\lnot p$ são idênticos, pois $\Diamond$ é definido
  como~$\lnot\Box\lnot$.}
\end{proof}

Na !!{derivation} acima, a etapa final ``$p$ por $\lnot\lnot p$''
abrevia
  \begin{derivation}
    & $\Log{K} \Proves \lnot \Box \lnot\lnot p \lif \Diamond \lnot
    p$\\
    & $\Log{K} \Proves \lnot\lnot p \liff p$ & \Taut\\
    & $\Log{K} \Proves \lnot\Box p \lif \Diamond \lnot p$ & por \olref{prop:rewriting}
  \end{derivation}
Os papéis de $!C(q)$, $!A$ e $!B$ em \olref{prop:rewriting} são
desempenhados aqui, respectivamente, por $\lnot \Box q \lif \Diamond
\lnot p$, $\lnot\lnot p$ e~$p$.

Quando !!a{formula} contém uma sub-!!{formula} $\lnot\Diamond !A$,
podemos substituí-la por $\Box\lnot !A$ usando
\olref{prop:rewriting}, pois $\Log{K} \Proves \lnot\Diamond !A \liff
\Box\lnot !A$. Indicaremos esta e outras substituições semelhantes
simplesmente por ``$\Box\lnot$ por $\lnot\Diamond$.''

A proposição seguinte justifica que podemos estabelecer resultados de
!!{derivability} esquematicamente. P.ex., a proposição anterior
não estabelece apenas que $\Log{K} \Proves \lnot\Box p \lif \Diamond
\lnot p$, mas que $\Log{K} \Proves \lnot\Box !A \lif \Diamond \lnot
!A$ para qualquer~$!A$.

\begin{prop}
  Se $!A$ é uma instância de substituição de $!B$ e
  $\Log{K} \Proves !B$, então $\Log{K} \Proves !A$.
\end{prop}

\begin{proof}
  É trabalhoso, mas rotineiro, verificar (por indução no comprimento
  da !!{derivation} de~$!B$) que aplicar uma substituição a uma
  !!{derivation} inteira também produz uma !!{derivation} correta.
  Mais precisamente, instâncias de substituição de instâncias
  tautológicas são elas próprias instâncias tautológicas; instâncias
  de substituição de instâncias de~\iftag{prvDiamond}{\Dual{} e~}{}\Ax{K}
  são elas próprias instâncias de~\iftag{prvDiamond}{\Dual{}
    e~}{}\Ax{K}; e as aplicações de \MP{} e~\Nec{} permanecem corretas
  quando !!{formula}s substituem !!{propositional variable}s tanto
  na(s) premissa(s) quanto na conclusão.
\end{proof}


\end{document}
